\chapter{A Tick Store on Open Formats}\label{ch:pl:a-tick-store-on-open-formats}

At 16:05 the day's capture files close. By 16:40 they must be compacted, sorted and queryable, because the overnight research
starts at 17:00, and a researcher who asks for the quote prevailing before each of this week's trades must not notice which
days come from the files written during the session and which from the compacted history. A week in the life of such a store
is this chapter: five simulated days go through capture, an \omterm{def:pl:a-tick-store-on-open-formats:split}{intraday store} written as the session runs, a compaction at the
close, a \omterm{def:pl:a-tick-store-on-open-formats:split}{historical store} of partitioned columnar files, and one query layer over both. On the fourth day the quotes gain a
field, and the store has to carry on. Nothing here needs a proprietary database: the store is built from the open formats of
chapter~3, an embedded query engine and a dataframe library, with a memory-mapped flat file for the hottest data, read the same
way from Python, C++ and Rust (\texttt{firm.tickstore}).

\section{Live capture: append-only files}

\begin{definition}[Tick store, append-only file]\label{def:pl:a-tick-store-on-open-formats:store}
A \emph{tick store}\index{tick store} is the system that keeps a firm's \omterm{def:pl:capturing-and-storing-tick-data:tick}{tick data} and answers queries on it, from the current
session to the oldest year retained. An \emph{append-only file}\index{append-only file} is a file whose writer only adds bytes at
its end and never changes what it has written, so that a reader can use every complete record already there while the writer
continues.
\end{definition}

During the session the store writes the way data arrives: in time order, as it comes, with nothing sorted, compressed or
indexed. The events of chapter~2 --- the receive time followed by the feed handler's 48-byte event --- go to a flat file of
fixed-width records; the quotes and trades derived from them (the top of book after each change, from Book~13's book builder,
and each execution with its price) go to Arrow IPC streams, the in-memory columnar format of chapter~3 written as a sequence of
record batches.

\begin{definition}[Flat file format, schema version]\label{def:pl:a-tick-store-on-open-formats:flat}
A \emph{flat file format}\index{flat file format} stores records of one fixed width one after another, after a header that
states their layout, so that record $i$ is at a computable offset and nothing needs to be parsed to reach it. The header's
\emph{schema version}\index{schema version} identifies the layout; a reader refuses a version it does not know.
\end{definition}

The store's flat file (\cref{fig:pl:a-tick-store-on-open-formats:flat}) starts with a magic number, the version, the record
width and the layout as JSON, padded to 64 bytes so that every record that follows is aligned the same way; version~1 records are
56 bytes, version~2 appends a venue, flags and padding for 64. \Cref{lst:pl:a-tick-store-on-open-formats:flat} is the writer and
the reader.

\omcode{code/firm/tickstore/firm_tickstore.py}{55}{83}{The flat file: a versioned header padded to 64 bytes, records appended
at the end, and a reader that maps the records instead of reading them.}\label{lst:pl:a-tick-store-on-open-formats:flat}

\begin{omfigure}
\begin{tikzpicture}[font=\scriptsize,f/.style={draw,minimum height=0.55cm,inner sep=2pt,align=center}]
\node[f,fill=gray!15,minimum width=1.0cm] (m) at (0.5,0) {OQTS};
\node[f,fill=gray!15,minimum width=1.3cm,anchor=west] (v) at (m.east) {v, width};
\node[f,fill=gray!10,minimum width=3.1cm,anchor=west] (s) at (v.east) {schema (JSON), pad};
\foreach \i/\c in {0/omDef,1/omMeth,2/omDef,3/omMeth,4/omDef} {\node[f,fill=\c!18,minimum width=1.2cm,anchor=west] at (5.4+\i*1.2,0) {rec \i};}
\node at (11.8,0) {\dots};
\draw[decorate,decoration={brace,amplitude=4pt}] (0,0.35) -- (5.4,0.35) node[midway,above=4pt]{header: multiple of 64 bytes};
\draw[decorate,decoration={brace,amplitude=4pt}] (5.4,0.35) -- (11.5,0.35) node[midway,above=4pt]{56 or 64 bytes each, appended};
\foreach \x in {5.4,9.4} {\draw[dashed,omThm] (\x,-0.35) -- (\x,-0.9);}
\node[omThm,anchor=north] at (7.4,-0.9) {one page of the mapping, read on first touch};
\draw[omThm,<->] (5.4,-0.75) -- (9.4,-0.75);
\end{tikzpicture}

\omcaption{The \omterm{def:pl:a-tick-store-on-open-formats:store}{tick store}'s flat file: a header of multiple of 64 bytes (magic, version, record width, the layout as JSON),
then fixed-width records appended as they arrive. A reader maps the file; the operating system reads a page from disk the first
time a record in it is touched.}\label{fig:pl:a-tick-store-on-open-formats:flat}
\end{omfigure}

\begin{remark}[Why two intraday forms]\label{rem:pl:a-tick-store-on-open-formats:twoforms}
The flat file is the fastest thing to write and to scan, and the layout the trading path already uses; the IPC stream carries a
schema, nullable and string columns, and is read into any Arrow-speaking tool without conversion. A store needs both because
its intraday readers differ: a C++ replay tool wants the flat records, a researcher's notebook wants a table. Each is written
once; nothing is converted during the session.
\end{remark}

\section{End-of-day compaction}

\begin{definition}[End-of-day compaction, sort key]\label{def:pl:a-tick-store-on-open-formats:compaction}
\emph{End-of-day compaction}\index{end-of-day compaction} rewrites a day's intraday files into the historical layout: partitioned,
sorted, encoded and compressed for reading. The \emph{sort key}\index{sort key} of a file is the list of columns its rows are
ordered by; the \omterm{def:pl:columnar-formats:pushdown}{zone maps} of chapter~3 prune well on its leading columns.
\end{definition}

The compaction of \texttt{firm.tickstore} (\cref{lst:pl:a-tick-store-on-open-formats:compact}) sorts each table by symbol and
time and writes one Parquet file per partition, under \texttt{<table>/date=YYYY-MM-DD/symbol=S/}: the directory names carry
the partition values, a convention the query engines read as columns. It is \cref{met:pl:columnar-formats:layout} of
chapter~3 made concrete: partition by date and symbol, sort by time inside.

\omcode{code/firm/tickstore/firm_tickstore.py}{147}{161}{End-of-day compaction: sort by symbol then time, and write one Parquet
file per date and symbol, the partition values in the directory names.}\label{lst:pl:a-tick-store-on-open-formats:compact}

\Cref{fig:pl:a-tick-store-on-open-formats:sizes} counts the bytes of the simulated week in each form. The busiest day, Monday,
has 258\,535 events; its intraday files take 17.5~MB (the flat events and the two IPC streams), its compacted partitions 6.15~MB,
2.85 times less, with the events at 20.0 bytes each in their partitions. Compaction of that day took 0.12~s on the laptop:
about 2.2 million events a second on one core, sorting, \omterm{def:pl:capturing-and-storing-tick-data:partition}{partitioning}, encoding and compressing included.

\begin{omfigure}
\begin{tikzpicture}
\begin{axis}[width=11.5cm,height=5.8cm,ybar,bar width=8pt,ylabel={MB},xlabel={trading day (month-day)},
  xtick={0,1,2,3,4},xticklabels from table={figdata/platforms/04-a-tick-store-on-open-formats/sizes.csv}{day},
  tick label style={font=\small},label style={font=\small},ymin=0,ymax=16,grid=major,grid style={gray!20},
  legend columns=3,legend style={font=\footnotesize,at={(0.5,-0.3)},anchor=north,draw=none,/tikz/every even column/.append style={column sep=8pt}},
  table/col sep=comma]
\addplot[fill=omDef!45,draw=omDef] table[x=k,y=flat]{figdata/platforms/04-a-tick-store-on-open-formats/sizes.csv};
\addplot[fill=omMeth!45,draw=omMeth] table[x=k,y=ipc]{figdata/platforms/04-a-tick-store-on-open-formats/sizes.csv};
\addplot[fill=omThm!45,draw=omThm] table[x=k,y=parquet]{figdata/platforms/04-a-tick-store-on-open-formats/sizes.csv};
\legend{flat events (intraday),IPC quotes and trades (intraday),Parquet partitions (historical)}
\end{axis}
\end{tikzpicture}

\omcaption{Bytes of each day of the simulated week in the intraday forms and, for the four compacted days, in the historical
Parquet partitions (events, quotes and trades together). The fifth day is the current one and has no historical form yet; its
flat records are 64 bytes, the version~2 layout. Data: \texttt{fig\_tickstore.py}.}\label{fig:pl:a-tick-store-on-open-formats:sizes}
\end{omfigure}

\begin{method}[Safe compaction]\label{met:pl:a-tick-store-on-open-formats:safe}
(1) Write the day's partitions to new files in a staging directory, never over the intraday files. (2) Verify them against
the intraday data: the same row count per table and symbol, the same range of sequence numbers with the same gaps, the same
sums of quantities. (3) Make them visible in one atomic step --- a rename of the staging directory, or a new version of a
manifest that lists the day's files. (4) Only then delete the intraday files, or keep them for the retention of the \omterm{def:pl:capturing-and-storing-tick-data:raw}{raw
capture} (chapter~2).
\end{method}

A reader that lists a directory while the compactor writes into it sees half a day; a reader of a manifest sees either the old
day or the new one. The verification step is cheap because the intraday files are still there, and it is the last moment at
which a lost partition can be rebuilt without going back to the \omterm{def:pl:capturing-and-storing-tick-data:raw}{raw capture}; chapter~7 turns such checks into the firm's
data-quality rules.

\begin{dated}{2026-09}{The proprietary design this one mirrors}\label{dat:pl:a-tick-store-on-open-formats:kdb}
The intraday/historical split is the design of the kdb+tick architecture of the proprietary kdb+ database, long common in banks
and trading firms: in its documentation (consulted September 2026) a tickerplant writes incoming records to a log and pushes
them to a real-time database, which holds the current day in memory and writes it to the historical database at the end of
day; the historical database gives access to all earlier days. This book teaches the open-format version of the same design.
\end{dated}

\section{The intraday and the historical store}

\begin{definition}[Intraday store, historical store]\label{def:pl:a-tick-store-on-open-formats:split}
The \emph{intraday store}\index{intraday store} holds the current session's data in the forms written during it, optimised for
appending and for recent reads; the \emph{historical store}\index{historical store} holds every compacted day, optimised for
reading ranges across days and symbols.
\end{definition}

The split is invisible to a reader only if one query layer covers both. \texttt{Store} opens an embedded DuckDB and creates one
view per table: the historical Parquet partitions, read with their directory names as columns, and the current day's IPC tables,
registered in memory, joined by a union that matches columns by name (\cref{lst:pl:a-tick-store-on-open-formats:store}). A
query on \texttt{quotes} then covers Monday to Friday, whichever store holds each day.

\omcode{code/firm/tickstore/firm_tickstore.py}{165}{187}{The query layer: one view per table over the historical partitions and
the current day's intraday table, matched by column name.}\label{lst:pl:a-tick-store-on-open-formats:store}

\begin{omfigure}
\begin{tikzpicture}[font=\scriptsize,>=Stealth,box/.style={draw,rounded corners=2pt,fill=omDef!10,minimum height=0.7cm,inner sep=3pt,align=center}]
\node[box] (feed) at (0,0.1) {feed lines\\(capture)};
\node[box,fill=omMeth!15] (flat) at (3.3,0.7) {flat events\\(append-only)};
\node[box,fill=omMeth!15] (ipc) at (3.3,-0.5) {IPC quotes, trades\\(record batches)};
\draw[dashed,gray] (1.9,-1.15) rectangle (4.7,1.35);
\node[font=\scriptsize\itshape,anchor=south] at (3.3,1.35) {intraday store};
\node[box,fill=omThm!12] (cmp) at (3.3,-2.4) {end-of-day\\compaction};
\node[box,fill=omProp!15,text width=2.6cm] (hist) at (7.2,-2.4) {Parquet\\date=D/symbol=S\\sorted by time};
\node[font=\scriptsize\itshape,anchor=north] at (7.2,-3.15) {historical store};
\node[box] (q) at (10.4,0.1) {query layer\\(DuckDB, Polars)};
\node[box,fill=gray!10] (u) at (10.4,1.5) {research, replay};
\draw[->] (feed) -- (flat);
\draw[->] (feed) -- (ipc);
\draw[->] (3.3,-1.15) -- (cmp);
\draw[->] (cmp) -- (hist);
\draw[->] (hist.east) -| (q.south);
\draw[->] (4.7,0.1) -- node[above]{today} (q.west);
\draw[->] (q) -- (u);
\end{tikzpicture}

\omcaption{The \omterm{def:pl:a-tick-store-on-open-formats:store}{tick store}: capture writes the \omterm{def:pl:a-tick-store-on-open-formats:split}{intraday store} as the session runs; compaction at the close turns each day into
sorted, partitioned columnar files; one query layer reads the historical partitions and the current day together.}\label{fig:pl:a-tick-store-on-open-formats:arch}
\end{omfigure}

\begin{example}[One window, two stores]\label{ex:pl:a-tick-store-on-open-formats:window}
The count of SIM2's quote changes from 09:32 to 09:35 on each day of the week comes back as 12\,291, 996, 1\,388, 1\,197 and
4\,473: four days from the Parquet partitions, the fifth from the intraday IPC table, in one SQL statement that names neither.
The query took 18~ms on the laptop.
\end{example}

\section{As-of joins at scale}

The as-of join of One Quant Book~7, chapter~3 --- each row of one table with the last row of another known at its time --- is
the \omterm{def:pl:a-tick-store-on-open-formats:store}{tick store}'s most frequent query: each trade with the quote that prevailed when it happened, for effective spreads, trade
signs and mark-outs. The store computes it three ways, which must agree row for row: DuckDB's \texttt{ASOF JOIN}, Polars'
\texttt{join\_asof}, and Book~7's \texttt{firm.pit.asof\_join} applied to each day and symbol.

\omcode{code/firm/tickstore/firm_tickstore.py}{190}{202}{The as-of join in SQL: each trade with the last quote of its day and
symbol of an earlier sequence number (key \texttt{seq}), or at or before its timestamp (key \texttt{ts}).}\label{lst:pl:a-tick-store-on-open-formats:asof}

The first version of the chapter joined on the timestamp, and the three engines disagreed. The reason is in the data: an
execution and the change of quote it causes carry the same exchange timestamp, so ``the last quote at or before the trade's
time'' can be the quote the trade itself produced, and which of two quotes with equal timestamps an engine returns is not
specified. The quote that prevailed \emph{before} the trade is the last one of an earlier sequence number, and the feed's
sequence number totally orders events.

\begin{proposition}[Join on the sequence number]\label{prop:pl:a-tick-store-on-open-formats:seq}
If every quote and trade row carries the sequence number of the event that produced it, the strict as-of join on sequence
number returns, for each trade, the top of book immediately before the event that executed it, and the result does not depend
on the join algorithm.
\end{proposition}

\begin{proof}
Sequence numbers are distinct across events and increase in the order the book processed them. The quote rows of a symbol with
sequence numbers below the trade's are the states of its book before the execution; the last of them is the state just before
it. Distinct keys leave no tie for an algorithm to break.
\end{proof}

On the week's 17\,281 trades, the three engines agree on every row with the sequence-number join. The timestamp join returns a
different quote for 223 of them, 1.3\%, each time the quote the execution itself had just produced: exactly the trades whose
effective spread a researcher would measure as zero or negative. On this data DuckDB's join took 16~ms, Polars' 45~ms, and the
Python reference 0.17~s (\cref{fig:pl:a-tick-store-on-open-formats:times}).

\begin{omfigure}
\begin{tikzpicture}
\begin{axis}[width=11.5cm,height=6.2cm,xbar,bar width=8pt,xmode=log,log origin=infty,xmin=0.0005,xmax=1,
  xlabel={seconds, best of several runs (log scale)},y dir=reverse,
  ytick={0,1,2,3,4,5,6},yticklabels from table={figdata/platforms/04-a-tick-store-on-open-formats/measured_tickstore.csv}{task},
  tick label style={font=\small},label style={font=\small},grid=major,grid style={gray!20},
  point meta=explicit,nodes near coords,nodes near coords style={font=\footnotesize,anchor=west,/pgf/number format/fixed,/pgf/number format/precision=4},
  table/col sep=comma]
\addplot[fill=omDef!45,draw=omDef] table[x=seconds,y expr=\coordindex,meta=seconds]{figdata/platforms/04-a-tick-store-on-open-formats/measured_tickstore.csv};
\end{axis}
\end{tikzpicture}

\omcaption{Time of the store's main tasks on the simulated week: the compaction of the busiest day (258\,535 events), the
five-day window query, the as-of join of 17\,281 trades by three engines, and a scan of one column of a day from the flat file
and from Parquet. Measured on a laptop (Intel Core Ultra 7 155H) under WSL2, one thread, machine otherwise idle. Data:
\texttt{bench\_tickstore.py}.}\label{fig:pl:a-tick-store-on-open-formats:times}
\end{omfigure}

\section{Flat files and memory mapping}

\begin{definition}[Memory-mapped file]\label{def:pl:a-tick-store-on-open-formats:mmap}
A \emph{memory-mapped file}\index{memory-mapped file} is a file mapped into a process's virtual address space, so that its bytes
are read as memory: the operating system loads each page from disk (or from its cache) the first time it is touched, and
processes that map the same file share one copy of its pages.
\end{definition}

A fixed-width file and a memory map fit together: record $i$ is at a known address, nothing is parsed, and only the pages a
query touches are read. Book~12's memory-mapped datasets (chapter~23) use the same mechanism for training data. The store's reader
in Python is \texttt{numpy.memmap}; the C++20 reader (\cref{lst:pl:a-tick-store-on-open-formats:cpp}) calls \texttt{mmap}
itself, and the Rust twin declares \texttt{mmap} and \texttt{munmap} from the C library, since the series' Rust uses no external
crates. The three compute the same per-instrument summary --- records, quantity executed, last sequence number --- on the same
fixtures of both versions. Scanning one column of Monday's events took 0.8~ms from the mapped flat file and 2.1~ms from its
Parquet partitions: the flat file is more than twice as fast to scan and 2.8 times larger to keep.

\omcode{code/firm/tickstore/cpp/firm_tickstore.hpp}{29}{62}{The C++20 reader: map the file, check the header, and read a record
at a computed offset; the page under it is loaded by the operating system on first touch.}\label{lst:pl:a-tick-store-on-open-formats:cpp}

\section{Schema evolution}

On Thursday the venue added a field to its quotes, and the store's quotes carry it from that day on. Monday to Wednesday have no
such column.

\begin{definition}[Schema evolution]\label{def:pl:a-tick-store-on-open-formats:evolution}
\emph{Schema evolution}\index{schema evolution} is the change of a dataset's layout over time under rules that let every reader
read every version: here, fields may be added at the end with a default, and are never removed, renamed, retyped or reordered.
\end{definition}

Under that rule a reader of the new version reads an old file by filling the new fields with their defaults, and a reader of the
old version reads a new file by ignoring the fields it does not know; the flat file does both by its version number
(\texttt{upgrade}), the query layer by matching columns by name (a missing column reads as null). The same week's window query
counts the venue on none of the first three days and on every quote of the last two. The registry's check
(\texttt{check\_evolution}) refuses the changes that would break old readers: a field retyped from 32 to 64 bits, a field
dropped, two fields swapped, a field inserted before existing ones.

\begin{method}[Evolving a tick store's schema]\label{met:pl:a-tick-store-on-open-formats:evolve}
(1) Add the field at the end, with a default that means ``not known''. (2) Increase the version in the header and register the new
layout next to the old one. (3) Write every reader to accept every registered version. (4) Never rewrite history to the new
version unless the old readers are retired; if a correction is needed, write the corrected days as new versions (chapter~7).
\end{method}

\section{Tutorial: a week in the store}\label{tut:pl:a-tick-store-on-open-formats:tut}

\begin{tutorial}
\textbf{Goal.} Run five simulated days through capture, the \omterm{def:pl:a-tick-store-on-open-formats:split}{intraday store}, compaction and the query layer, with a schema change
on the fourth day, and check the as-of join three ways. \textbf{End state:} \cref{fig:pl:a-tick-store-on-open-formats:sizes},
the window of \cref{ex:pl:a-tick-store-on-open-formats:window}, and three identical joins of 17\,281 trades.
\begin{tutsteps}
\item \textbf{Build the week}: \texttt{pl\_tickstore.build(root)} captures each day (\texttt{capture\_day}: events to a flat
file, quotes and trades to IPC streams, one flush every 20\,000 events) and compacts the first four
(\cref{lst:pl:a-tick-store-on-open-formats:compact}).
\item \textbf{Query}: \texttt{window(root)} counts SIM2's quotes in the same three minutes of each day.
\item \textbf{Join}: \texttt{asof\_all(root)} runs DuckDB, Polars and the reference; then rerun DuckDB with the timestamp as the key
and count the trades whose quote changes.
\item \textbf{Map}: build and run \texttt{firm\_tickstore\_test.cpp} and \texttt{cargo test} in \texttt{code/firm/tickstore/rust}
on the fixtures; compare with \texttt{flat\_summary}.
\item \textbf{Measure}: \texttt{bench\_tickstore.py} times the compaction and the queries.
\end{tutsteps}
\textbf{What to change next.} Add a third \omterm{def:pl:a-tick-store-on-open-formats:flat}{schema version} with a condition code and read the whole week; make the compaction
write two sort orders (symbol then time, and time then symbol) and compare a cross-sectional query on each.
\end{tutorial}

\section{Build: the tick store}\label{bld:pl:a-tick-store-on-open-formats:tickstore}

\begin{build}
\bfield{Purpose} The firm's \omterm{def:pl:a-tick-store-on-open-formats:store}{tick store}: every later chapter of this book that reads market data reads it here, the backtest
engine of chapter~11 through its data-access layer, the quality rules of chapter~7 on its partitions.
\bfield{Interface} \texttt{EVENT\_V1}, \texttt{EVENT\_V2}, \texttt{SCHEMAS}; \texttt{FlatWriter(path, version)},
\texttt{flat\_open(path) -> (version, memmap)}, \texttt{upgrade}; \texttt{IpcWriter}, \texttt{ipc\_read};
\texttt{quotes\_trades(events, symbols)}; \texttt{compact(tables, root, date, row\_group\_size)}; \texttt{Store(root, intraday,
date).sql(q)}; \texttt{asof\_duckdb}, \texttt{asof\_polars}, \texttt{asof\_reference}; \texttt{check\_evolution}. C++20
\texttt{firm::tickstore::FlatFile}, \texttt{summarise}; Rust \texttt{FlatFile::open}, \texttt{get}, \texttt{summarise}.
\bfield{Rules} Intraday files are append-only; the header states the version and a reader refuses an unknown one; historical
files are written once per day and partition; every quote and trade carries its event's sequence number; the as-of join is
strict on the sequence number; schemas only grow at the end.
\bfield{Acceptance tests} \texttt{code/firm/tickstore/tests/}, \texttt{cpp/firm\_tickstore\_test.cpp}, \texttt{rust/}: flat files of
both versions written, mapped and refused when foreign; the Python reader equal to the fixture the twins check; quotes and trades
from a hand-made book; an IPC round trip; the query layer over two historical days of different schemas and the current day;
three identical joins and the timestamp join's tie; the evolution rules.
\bfield{Stretch} A second sort order written by the compaction; a sorted merge of intraday and historical rows without a union;
compaction of a day while the next day's capture runs.
\end{build}

\begin{omsources}
\item KX, \emph{kdb+tick} documentation (tickerplant, real-time and historical databases).
\item DuckDB, \emph{AsOf Join} (documentation); Polars, \texttt{DataFrame.join\_asof} (API reference).
\item Apache Arrow, \emph{IPC streaming format} (specification); \texttt{mmap(2)}, Linux manual page.
\end{omsources}

\section{Exercises}

\begin{exercise}[$\star$]\label{exo:pl:a-tick-store-on-open-formats:1}
Monday's flat file is 14\,478\,152 bytes for 258\,535 records of version~1. How long is its header, and how long would the file
be in version~2?
\end{exercise}

\begin{exercise}[$\star$]\label{exo:pl:a-tick-store-on-open-formats:2}
Why may a reader use an \omterm{def:pl:a-tick-store-on-open-formats:store}{append-only file} while the writer is still writing it, and what must it do with a partial last record?
\end{exercise}

\begin{exercise}[$\star$]\label{exo:pl:a-tick-store-on-open-formats:3}
Monday's intraday files take 17.5~MB and its compacted partitions 6.15~MB. What is the ratio, and where does the saving come
from?
\end{exercise}

\begin{exercise}[$\star\star$]\label{exo:pl:a-tick-store-on-open-formats:4}
A researcher computes effective spreads with an as-of join on the timestamp. On the week's data, how many trades get the wrong
quote, and in which direction are their spreads biased?
\end{exercise}

\begin{exercise}[$\star\star$]\label{exo:pl:a-tick-store-on-open-formats:5}
At the measured compaction rate, how many events can the store compact between 16:05 and 16:40 on one core, and how does that
compare with the options feed's planned 311 billion messages a day of chapter~2?
\end{exercise}

\begin{exercise}[$\star\star$]\label{exo:pl:a-tick-store-on-open-formats:6}
A new field is inserted in the middle of the record instead of appended. What does an old reader see, and what does
\texttt{check\_evolution} report?
\end{exercise}

\begin{exercise}[$\star\star\star$]\label{exo:pl:a-tick-store-on-open-formats:7}
\emph{Coding.} With \texttt{flat\_summary} and the C++ reader, check the current day's flat file: how many records does each
instrument have, and what quantity was executed on instrument~2?
\end{exercise}

\begin{exercise}[$\star\star\star$]\label{exo:pl:a-tick-store-on-open-formats:8}
\emph{Find the flaw.} ``To save space we compact the day in place: the compactor reads the intraday files, writes the Parquet
partitions over them, and a query during compaction reads whatever is there.''
\end{exercise}

\section{Problem: The 16:40 Deadline}

\begin{problem}[{Weekend problem --- a week in the tick store}]\label{pb:pl:a-tick-store-on-open-formats:1}
The chapter's simulated week, its store and the measurements of \texttt{bench\_tickstore.py}.

\medskip\noindent\textbf{Part I --- Intraday.}
\begin{enumerate}
\item How many events, quotes and trades does the week hold?
\item What is in a flat file's header, and why is it padded to 64 bytes?
\item Why does the \omterm{def:pl:a-tick-store-on-open-formats:split}{intraday store} write two forms, and which readers use which?
\item How many bytes is a record in version~1 and in version~2?
\item Why is nothing sorted or compressed during the session?
\end{enumerate}

\medskip\noindent\textbf{Part II --- Compaction.}
\begin{enumerate}[resume]
\item What \omterm{def:pl:a-tick-store-on-open-formats:compaction}{sort key} and \omterm{def:pl:capturing-and-storing-tick-data:partition}{partition key} does the compaction use, and why those?
\item How many bytes does Monday take before and after compaction?
\item How long did Monday's compaction take on the laptop, and at what rate?
\item How many events could the 35 minutes from 16:05 to 16:40 absorb at that rate on one core?
\item What would you change to compact a full options feed in the window?
\end{enumerate}

\medskip\noindent\textbf{Part III --- Queries.}
\begin{enumerate}[resume]
\item How does one SQL statement read both stores?
\item How many quote changes does SIM2 have from 09:32 to 09:35 each day?
\item Why must the as-of join use the sequence number rather than the timestamp?
\item How many trades change quote between the two joins, and which engines agree?
\item How long did the three joins take?
\end{enumerate}

\medskip\noindent\textbf{Part IV --- The verdict.}
\begin{enumerate}[resume]
\item State the \emph{named result}: the compaction rate and the size reduction of the busiest day, the query latencies, and the
largest day the window absorbs.
\item How do old and new quotes coexist after Thursday's schema change?
\item What does the memory-mapped flat file buy over Parquet for the current day, and what does it cost?
\item Which part of this design comes from the proprietary tick databases, and which from the open formats?
\item In one sentence: what makes a \omterm{def:pl:a-tick-store-on-open-formats:store}{tick store}'s two halves look like one?
\end{enumerate}
\end{problem}

\section{Interview questions}

\begin{interviewq}[$\star$ \iqroles{developer, researcher}]\label{iq:pl:a-tick-store-on-open-formats:1}
Why do \omterm{def:pl:a-tick-store-on-open-formats:store}{tick stores} separate the current day from history?
\end{interviewq}

\begin{interviewq}[$\star\star$ \iqroles{developer, researcher}]\label{iq:pl:a-tick-store-on-open-formats:2}
You join trades to quotes as of the trade time and some effective spreads come out negative. What is wrong?
\end{interviewq}

\begin{interviewq}[$\star\star$ \iqroles{developer}]\label{iq:pl:a-tick-store-on-open-formats:3}
What does memory-mapping a file give you, and when is it a bad idea?
\end{interviewq}

\begin{interviewq}[$\star\star$ \iqroles{developer}]\label{iq:pl:a-tick-store-on-open-formats:4}
A venue adds a field to its feed. How do you change the store without breaking five years of readers?
\end{interviewq}

\begin{interviewq}[$\star\star$ \iqroles{developer}]\label{iq:pl:a-tick-store-on-open-formats:5}
How would you partition and sort a year of quotes for 8\,000 symbols, and why?
\end{interviewq}

\begin{interviewq}[$\star\star\star$ \iqroles{developer}]\label{iq:pl:a-tick-store-on-open-formats:6}
Design a \omterm{def:pl:a-tick-store-on-open-formats:store}{tick store} on open formats for a firm without a proprietary database: capture, intraday, compaction, historical layout,
query layer, and what you would measure first.
\end{interviewq}
